%%%%%%%%%%%%%%%%%%%%%%%%%%%%%%%%%%%%%%%%%%%%%%%%%%%%%%%%%%%%%%
%% Chapter 8: Conclusions
%%%%%%%%%%%%%%%%%%%%%%%%%%%%%%%%%%%%%%%%%%%%%%%%%%%%%%%%%%%%%%

\section{Conclusions}
\label{chap:conclusions}

\subsection{Answers to thesis goals}
\label{sec:conclusions-goals}

Section~\ref{sec:goals-scope} set out five goals for this project. Two of them
are addressed in full by the work presented in this thesis, and three remain
open at the time of writing.

The annotation platform goal is met by the system described in
Chapter~\ref{chap:annotation-platform}: a working, decentralized data
collection application supporting dataset, video, gesture class, and gesture
type management, asynchronous MediaPipe-based landmark extraction, and a
browser-based clip annotation editor, all deployable locally through Docker
Compose.

The isolated video-to-gloss recognition goal is met by the system described in
Chapter~\ref{chap:isolated-gloss-recognition}: a retrieval-based, multi-objective
architecture benchmarked across five temporal encoders on a 32-gloss subset of
the ASL Citizen dataset, with the Transformer encoder achieving 90.75\% 1-NN
accuracy and per-class recall exceeding 0.75 across all evaluated glosses.

\todo{PLACEHOLDER: Once Chapters~\ref{chap:text-translation},
\ref{chap:sentence-level-translation}, and~\ref{chap:educational-platform} are
completed, summarise here how (and how well) each of the remaining goals --
gloss-text translation in both directions, end-to-end sentence-level
translation, and the educational demonstration platform -- was met, with
reference to their respective results sections.}

\subsection{Limitations}
\label{sec:conclusions-limitations}

The isolated gloss recognition system's principal limitation, identified
through confusion-matrix analysis, per-class recall, t-SNE visualisation, and
positive/negative pair distance tracking
(Section~\ref{sec:gloss-recognition-discussion}), is intra-class variability
arising from signer-dependent execution differences: unseen signers produce
more dispersed embeddings than seen signers, even though the overall cluster
structure generalises. The benchmark is also currently limited to a 32-gloss
vocabulary subset of ASL Citizen, and to isolated, pre-segmented gestures rather
than continuous signing.

The annotation platform's known limitations
(Section~\ref{sec:annotation-system-architecture}) include incomplete dataset
export processing, video metadata (FPS, duration) not yet fully derived from
uploaded files, no authentication or authorisation, limited validation around
clip ranges and dataset consistency, and limited automated test coverage.

\todo{PLACEHOLDER: Once Chapters~\ref{chap:text-translation},
\ref{chap:sentence-level-translation}, and~\ref{chap:educational-platform} are
completed, add their specific limitations here (e.g.\ gloss-text corpus size
and coverage, end-to-end error propagation between recognition and translation
stages, educational platform usability constraints).}

\subsection{Future work}
\label{sec:conclusions-future-work}

Several concrete extensions to the isolated gloss recognition system are
identified in Section~\ref{sec:gloss-recognition-discussion}: replacing or
augmenting the KNN classifier with Dynamic Time Warping (DTW)-based sequence
alignment at retrieval time; replacing the MSE reconstruction loss with
soft-DTW as a differentiable training objective; adopting FAISS for scalable
approximate nearest-neighbour retrieval beyond the current 32-gloss benchmark;
introducing Graph Convolutional Networks to encode the skeletal graph topology
explicitly rather than relying on flattened landmark vectors; evaluating
quaternion-valued LSTMs for rotation-aware temporal encoding; and extending the
system from isolated to continuous recognition via CTC, forming the basis of
the sentence-level pipeline in Chapter~\ref{chap:sentence-level-translation}.
Notably, a live webcam recognition demo combining the Transformer encoder with
FAISS retrieval has already been implemented in the project repository since
this thesis's research report was written, ahead of some of the above being
formally written up.

For the annotation platform, future work should complete the dataset export
processing pipeline, strengthen import/export job status tracking, derive video
metadata directly from uploaded files rather than using fixed values, add
authentication and authorisation, and broaden automated test coverage across
the API, worker, and frontend editor.

\todo{PLACEHOLDER: Once Chapters~\ref{chap:text-translation},
\ref{chap:sentence-level-translation}, and~\ref{chap:educational-platform} are
completed, add their future work items here.}
